\problema{Serialize and Deserialize Binary Tree}{297}{src/02-arboles/297-serialize-and-deserialize-binary-tree.cpp}{D}

El proceso de \emph{serializar} consiste en transformar una estructura de datos
(como un árbol) en un formato lineal (como una cadena de texto) que pueda ser
enviado por la red o guardado en disco. El proceso inverso,
\emph{deserializar}, toma esa cadena y reconstruye la estructura original en
memoria de forma exacta. Diseña el algoritmo para lograr esto con un árbol binario.

\paragraph{Intuición.} Existen varias formas válidas de recorrer un árbol y
plasmarlo en texto (BFS por niveles como lo hace LeetCode internamente, o DFS).
El enfoque más elegante y rápido de programar es el recorrido DFS en
\textbf{preorden} (Raíz, Izquierda, Derecha).

Si recorremos el árbol así, y guardamos un carácter especial (digamos \texttt{"N"})
cada vez que nos topemos con un puntero nulo, la cadena resultante tendrá toda
la información necesaria para reconstruirse a sí misma sin ambigüedades.

\begin{center}
\ttfamily
\begin{tabular}{ccc}
 & 1 & \\
2 & & 3\\
\end{tabular}
\end{center}

\noindent Usando preorden y comas como separadores, el árbol de arriba genera
la cadena: \texttt{"1,2,N,N,3,N,N,"}.

\paragraph{Solución.} Para la serialización, usamos una recursión simple en
preorden, concatenando el valor del nodo y las llamadas recursivas.

Para la deserialización, el truco está en \emph{consumir} la cadena en el mismo
orden (preorden). El primer valor que leamos será obligatoriamente la raíz.
Después de crear la raíz, llamamos a la recursión para que consuma de la cadena
todo lo que pertenezca al subárbol izquierdo, y luego lo que pertenezca al
subárbol derecho. Usar un \texttt{stringstream} de C++ (pasado por referencia)
hace que "consumir tokens separados por comas" sea extremadamente sencillo y
elegante.

\lstinputlisting{src/02-arboles/297-serialize-and-deserialize-binary-tree.cpp}

\complejidad{$O(n)$}{$O(n)$}

\paragraph{Notas de entrevista (Amazon).} Al ser un problema "Hard", los
entrevistadores buscan ver tu dominio de las herramientas del lenguaje. Usar
\texttt{std::stringstream} con \texttt{getline} en C++ demuestra madurez. Otro
punto a favor: menciona por qué el recorrido BFS (por niveles) también es
válido pero suele requerir más código (colas y manejo de nulos por nivel),
mientras que DFS preorden brilla por la brevedad de su implementación recursiva.
Evita usar variables globales enteras para el índice del string; en su lugar,
pasa el estado (el stream o un iterador) por referencia.
